
\documentclass{article}
\usepackage{colortbl}
\usepackage{makecell}
\usepackage{multirow}
\usepackage{supertabular}

\begin{document}

\newcounter{utterance}

\centering \large Interaction Transcript for game `hot\_air\_balloon', experiment `air\_balloon\_survival\_en\_complexity\_hard', episode 2 with Aya{-}Expanse{-}32B{-}t0.0.
\vspace{24pt}

{ \footnotesize  \setcounter{utterance}{1}
\setlength{\tabcolsep}{0pt}
\begin{supertabular}{c@{$\;$}|p{.15\linewidth}@{}p{.15\linewidth}p{.15\linewidth}p{.15\linewidth}p{.15\linewidth}p{.15\linewidth}}
    \# & $\;$A & \multicolumn{4}{c}{Game Master} & $\;\:$B\\
    \hline

    \theutterance \stepcounter{utterance}  
    & & \multicolumn{4}{p{0.6\linewidth}}{
        \cellcolor[rgb]{0.9,0.9,0.9}{
            \makecell[{{p{\linewidth}}}]{
                \texttt{\tiny{[P1$\langle$GM]}}
                \texttt{You are participating in a collaborative negotiation game.} \\
\\ 
\texttt{Together with another participant, you must agree on a single set of items that will be kept. Each of you has your own view of how much each item matters to you (importance). You do not know how the other participant values the items. Additionally, you are given the effort each item demands.} \\
\texttt{You may only agree on a set if the total effort of the selected items does not exceed a shared limit:} \\
\\ 
\texttt{LIMIT = 9322} \\
\\ 
\texttt{Here are the individual item effort values:} \\
\\ 
\texttt{Item effort = \{"A08": 380, "B95": 448, "A39": 568, "A65": 351, "B81": 896, "B34": 821, "A28": 526, "C82": 502, "B69": 326, "A87": 962, "B64": 619, "B15": 115, "C79": 598, "B16": 662, "C53": 300, "C83": 830, "C10": 562, "A97": 679, "A25": 283, "C84": 440, "A71": 12, "B23": 863, "B96": 319, "C51": 771, "C17": 89, "B19": 655, "C72": 504, "A94": 118, "A46": 514, "B43": 226, "B37": 892, "A88": 621, "A96": 766, "C60": 659, "A75": 768\}} \\
\\ 
\texttt{Here is your personal view on the importance of each item:} \\
\\ 
\texttt{Item importance values = \{"A08": 476, "B95": 544, "A39": 664, "A65": 447, "B81": 992, "B34": 917, "A28": 622, "C82": 598, "B69": 422, "A87": 1058, "B64": 715, "B15": 211, "C79": 694, "B16": 758, "C53": 396, "C83": 926, "C10": 658, "A97": 775, "A25": 379, "C84": 536, "A71": 108, "B23": 959, "B96": 415, "C51": 867, "C17": 185, "B19": 751, "C72": 600, "A94": 214, "A46": 610, "B43": 322, "B37": 988, "A88": 717, "A96": 862, "C60": 755, "A75": 864\}} \\
\\ 
\texttt{Goal:} \\
\\ 
\texttt{Your goal is to negotiate a shared set of items that benefits you as much as possible (i.e., maximizes total importance to YOU), while staying within the LIMIT. You are not required to make a PROPOSAL in every message {-} you can simply negotiate as well. All tactics are allowed!} \\
\\ 
\texttt{Interaction Protocol:} \\
\\ 
\texttt{You may only use the following structured formats in a message:} \\
\\ 
\texttt{PROPOSAL: \{'A', 'B', 'C', …\}} \\
\texttt{Propose keeping exactly those items.} \\
\texttt{REFUSE: \{'A', 'B', 'C', …\}} \\
\texttt{Explicitly reject opponent's proposal.} \\
\texttt{ARGUMENT: \{'...'\}} \\
\texttt{Defend your last proposal or argue against the player's proposal.} \\
\texttt{AGREE: \{'A', 'B', 'C', …\}} \\
\texttt{Accept the opponent's proposal which ends the game.} \\
\texttt{STRATEGIC REASONING: \{'...'\}} \\
\texttt{	Describe your strategic reasoning or anticipation explaining your choice of action. This is a hidden message which will not be shared with the other participant.} \\
\\ 
\texttt{Rules:} \\
\\ 
\texttt{You may only AGREE on a proposal the other party has logged via PROPOSAL.} \\
\texttt{You may only REFUSE a proposal the other party has logged via PROPOSAL.} \\
\texttt{Total effort of any PROPOSAL or AGREE set must be ≤ LIMIT.} \\
\texttt{Do NOT reveal your hidden importance scores.} \\
\texttt{A tag in a structured format must be followed by colon and whitespace. The argument must be a python set containing 0 or more strings.} \\
\texttt{So, it must be of the form TAG: \{...\}} \\
\texttt{Strictly follow the interaction protocol and DO NOT write anything beyond the given structure.} \\
\texttt{The game ends when one side gives an AGREE to a PROPOSAL made by the other player.} \\
\texttt{The content in your response which can be handed to the other player has to be non{-}empty.} \\
\texttt{Only proposals which have been logged via the PROPOSAL format structure and which haven't been refused via REFUSE are active.} \\
\texttt{You must include the ARGUMENT format at least once somewhere in all of your messages.} \\
\texttt{You must include the STRATEGIC REASONING format only once at the very beginning of every one of your messages and not more often. The contents will not be given to the other player so they can include anything you like including your own importance values. Here you should reason step by step to come up with you next move.} \\
            }
        }
    }
    & \\ \\

    \theutterance \stepcounter{utterance}  
    & \multicolumn{4}{p{0.6\linewidth}}{
        \cellcolor[rgb]{0.9,0.9,0.9}{
            \makecell[{{p{\linewidth}}}]{
                \texttt{\tiny{[P1$\rangle$GM]}}
                \texttt{**STRATEGIC REASONING:**} \\
\texttt{Given the item effort values and my personal importance scores, I need to strategize to maximize my benefits while staying within the effort limit. Initially, I'll focus on items with high importance scores and low effort values. I'll also need to anticipate the other participant's moves and adjust my strategy accordingly.} \\
\\ 
\texttt{**PROPOSAL:**} \\
\texttt{\{"A08", "A65", "B15", "B69", "C53", "C84", "A71", "C17"\}} \\
\\ 
\texttt{**ARGUMENT:**} \\
\texttt{\{"A08", "A65", "B15", "B69", "C53", "C84", "A71", "C17"\} are proposed as they have a combined effort of 3055, well below the limit, and a high combined importance score according to my valuation. This set includes items with low effort and high importance, providing a good starting point for negotiation.} \\
\\ 
\texttt{**STRATEGIC REASONING:**} \\
\texttt{This initial proposal sets a baseline. I can now gauge the other participant's response to determine their preferences and adjust my strategy. If they refuse, I'll look for items they might value highly and adjust my next proposal accordingly. If they agree, we end the game. If they propose an alternative, I'll evaluate it based on my importance scores and the effort limit.} \\
            }
        }
    }
    & & \\ \\

    \theutterance \stepcounter{utterance}  
    & & \multicolumn{4}{p{0.6\linewidth}}{
        \cellcolor[rgb]{0.9,0.9,0.9}{
            \makecell[{{p{\linewidth}}}]{
                \texttt{\tiny{[P1$\langle$GM]}}
                \texttt{Your response did not start with the proper strategic reasoning tag at the very beginning of your response. The very first structured format must be of the form STRATEGIC REASONING: \{...\}. Try again.} \\
            }
        }
    }
    & \\ \\

    \theutterance \stepcounter{utterance}  
    & \multicolumn{4}{p{0.6\linewidth}}{
        \cellcolor[rgb]{0.9,0.9,0.9}{
            \makecell[{{p{\linewidth}}}]{
                \texttt{\tiny{[P1$\rangle$GM]}}
                \texttt{**STRATEGIC REASONING:**} \\
\texttt{Given the item effort values and my personal importance scores, I need to strategize to maximize my benefits while staying within the effort limit of 9322. Initially, I'll focus on identifying items with high importance to me and relatively low effort requirements. This approach will allow me to propose a set of items that are likely to be beneficial to me while also leaving room for negotiation. I'll start by selecting items that have a combined effort well below the limit, ensuring that we have space to add items that the other participant may value highly.} \\
\\ 
\texttt{**PROPOSAL:**} \\
\texttt{\{"A08", "A65", "B15", "B69", "C53", "C84", "A71", "C17"\}} \\
\\ 
\texttt{**ARGUMENT:**} \\
\texttt{I propose keeping the items \{"A08", "A65", "B15", "B69", "C53", "C84", "A71", "C17"\} as they have a combined effort of 3055, which is well below the limit, and a high combined importance score according to my valuation. This set includes items with low effort and high importance, providing a good starting point for negotiation.} \\
\\ 
\texttt{By offering this set initially, I aim to demonstrate a willingness to compromise while also securing items that are valuable to me. I anticipate that the other participant may have different preferences, so I'm prepared to adjust my strategy based on their response. If they refuse this proposal, I'll look for items they might value highly and incorporate those into my next proposal. If they agree, we can end the game. If they make an alternative proposal, I'll evaluate it based on my importance scores and the effort limit.} \\
            }
        }
    }
    & & \\ \\

    \theutterance \stepcounter{utterance}  
    & & \multicolumn{4}{p{0.6\linewidth}}{
        \cellcolor[rgb]{0.9,0.9,0.9}{
            \makecell[{{p{\linewidth}}}]{
                \texttt{\tiny{[P1$\langle$GM]}}
                \texttt{Your response did not start with the proper strategic reasoning tag at the very beginning of your response. The very first structured format must be of the form STRATEGIC REASONING: \{...\}. Try again.} \\
            }
        }
    }
    & \\ \\

    \theutterance \stepcounter{utterance}  
    & \multicolumn{4}{p{0.6\linewidth}}{
        \cellcolor[rgb]{0.9,0.9,0.9}{
            \makecell[{{p{\linewidth}}}]{
                \texttt{\tiny{[P1$\rangle$GM]}}
                \texttt{Apologies for the oversight. Let's start afresh:} \\
\\ 
\texttt{**STRATEGIC REASONING:**} \\
\texttt{To initiate the negotiation, I will first propose a set of items that are highly important to me while ensuring the total effort remains significantly below the limit. This strategy allows room for flexibility and potential trade{-}offs during the negotiation process. I will select items with a combined effort value of approximately one{-}third of the total limit to account for potential adjustments.} \\
\\ 
\texttt{**PROPOSAL:**} \\
\texttt{\{"A08", "A65", "B15", "B69", "C53", "C84", "A71", "C17"\}} \\
\\ 
\texttt{**ARGUMENT:**} \\
\texttt{I propose the set \{"A08", "A65", "B15", "B69", "C53", "C84", "A71", "C17"\} as it meets my criteria of high importance and low effort. The combined effort of these items is 3055, leaving a substantial margin under the 9322 limit. This proposal serves as a starting point and demonstrates my intention to find a mutually beneficial agreement.} \\
\\ 
\texttt{I anticipate that the other participant may have different priorities, so I am prepared to adapt my strategy based on their response. If they refuse this proposal, I will analyze their counter{-}proposal or preferences and adjust my selection accordingly. My goal is to reach a consensus that maximizes the value for both parties within the given constraints.} \\
            }
        }
    }
    & & \\ \\

    \theutterance \stepcounter{utterance}  
    & & & \multicolumn{2}{p{0.3\linewidth}}{
        \cellcolor[rgb]{0.9,0.9,0.9}{
            \makecell[{{p{\linewidth}}}]{
                \texttt{\tiny{[GM$|$GM]}}
                \texttt{abort} \\
            }
        }
    }
    & & \\ \\

    \theutterance \stepcounter{utterance}  
    & & & \multicolumn{2}{p{0.3\linewidth}}{
        \cellcolor[rgb]{0.9,0.9,0.9}{
            \makecell[{{p{\linewidth}}}]{
                \texttt{\tiny{[GM$|$GM]}}
                \texttt{end game} \\
            }
        }
    }
    & & \\ \\

\end{supertabular}
}

\end{document}
